% Pembahasan
\section{Pembahasan}

\subsection{Interpretasi Nilai \texorpdfstring{$r \approx 0,300$}{r ≈ 0,300} sebagai Kemungkinan Batas Bawah}
Estimasi rerata korelasi \textit{off-diagonal} sebesar 0,300 tergolong moderat-rendah untuk delapan subtes yang mengukur kemampuan kognitif dan akademik yang saling berkaitan. Dugaan bahwa nilai ini tidak mencerminkan korelasi populasi riil bersumber dari bukti internal paper ini. Bukti itu adalah rerata sampel yang melampaui rerata nasional tidak resmi pada seluruh subtes dan $SD$ sampel yang menyusut pada mayoritas subtes (Tabel~\ref{tab:deskriptif}).

Rendahnya korelasi teramati itu paling rasional dijelaskan melalui mekanisme \textit{restriction of range}, sebagaimana ditunjukkan pada Tabel~\ref{tab:deskriptif}. Peserta yang memublikasikan skornya di media sosial pada dasarnya terseleksi berdasarkan skor komposit yang tinggi, dan skor komposit secara matematis merupakan fungsi dari kedelapan subtes. Seleksi sampel berdasarkan variabel yang berkorelasi dengan seluruh komponen pembentuknya menekan kovarians dan korelasi antarkomponen ke arah nol. Dengan demikian, nilai $r \approx 0,300$ paling tepat ditafsirkan sebagai batas bawah korelasi populasi riil.

Dukungan kuantitatif atas argumen batas bawah itu diperoleh melalui penyesuaian Thorndike Case III (Tabel~\ref{tab:korelasi_terkoreksi}). Koreksi atenuasi akibat \textit{restriction of range} pada skor komposit menaikkan rerata korelasi \textit{off-diagonal} sebesar 0,125 poin, sehingga korelasi \textit{raw sample} teratenuasi sebesar itu. Arah dan besaran koreksi ini konsisten dengan interpretasi bahwa rendahnya korelasi \textit{raw sample} lebih mencerminkan artefak bias seleksi mandiri daripada struktur psikometri UTBK-SNBT yang menyimpang dari kaidah tes kognitif pada umumnya. Namun, koreksi ini bertumpu pada asumsi homogenitas korelasi yang tidak terpenuhi sepenuhnya (lihat Subbagian Struktur Faktor dan Keterbatasan Homogenitas Korelasi).

\subsection{Dinamika Fluktuasi dan Pelandaian Korelasi Sepanjang Pengumpulan Data}
Pola pada Tabel~\ref{tab:saturasi_tahap} memperkuat konfirmasi empiris mengenai fenomena \textit{restriction of range}. Fluktuasi estimasi korelasi sepanjang tahap awal pengumpulan data per-PTN secara langsung merefleksikan komposisi heterogenitas skor yang masuk ke dalam sampel kumulatif. Penambahan kelompok skor dari institusi dengan profil selektivitas yang distingtif memperluas rentang skor komposit total sampel, mereduksi tingkat keparahan atenuasi, dan menaikkan korelasi teramati. Sebaliknya, akumulasi data dari kelompok dengan rentang kemampuan yang relatif homogen menekan korelasi ke arah bawah.

Namun, estimasi korelasi tidak lagi beranjak ketika penjaringan diperluas pada Tahap 18 dan Tahap 19 yang secara metodologis dirancang inklusif. Proporsi entri dengan skor komposit di bawah 550 tetap minim. Pola ini menunjukkan bahwa hambatan representativitas berakar pada struktur sosiologis audiens platform. Peserta yang aktif berbagi atau meninjau informasi seleksi masuk perguruan tinggi di media sosial telah tersaring secara sukarela pada kelompok berkinerja akademik menengah ke atas.

\subsection{Diferensiasi Kemampuan pada Skor Tinggi sebagai Penjelasan Pelengkap}
Selain \textit{restriction of range}, rendahnya korelasi pada sampel yang berisi skor tinggi dapat dijelaskan oleh mekanisme yang berbeda, yakni korelasi antartes yang tidak konstan di sepanjang rentang kemampuan. Gagasan itu adalah SLODR (lihat Metode). Korelasi antartes dan muatan \textit{g-factor} cenderung lebih rendah pada kelompok berkemampuan tinggi \autocite{spearman1927,detterman1989}. Meta-analisis \textcite{blum2017} mendukung pola itu, meskipun dukungannya tidak seragam di seluruh studi yang ditinjau.

Kedua mekanisme sama-sama memprediksi korelasi sampel yang rendah pada data ini, dengan perbedaan mendasar. \textit{Restriction of range} merupakan artefak seleksi yang tetap muncul walaupun korelasi populasi konstan, sedangkan SLODR berarti korelasi populasi itu sendiri lebih kecil pada skor tinggi. Karena seluruh sampel berasal dari wilayah skor tinggi (Tabel~\ref{tab:deskriptif}), data ini tidak memuat informasi untuk memisahkan keduanya. Oleh karena itu, SLODR berstatus penjelasan pelengkap yang tidak dapat dibedakan secara empiris dari bias seleksi.

Gagasan ini membawa dua implikasi bagi analisis di atas, yang diperiksa melalui simulasi skenario (Subbagian Sensitivitas terhadap Korelasi yang Bergantung pada Tingkat Skor). Implikasi pertama menyangkut koreksi Case III, yang mengasumsikan hubungan linear antara skor komposit sebagai variabel seleksi dan tiap subtes, dengan sebaran galat yang sama di seluruh rentang skor. Dengan asumsi itu, satu struktur korelasi dianggap berlaku di seluruh rentang skor. Pada model satu faktor sederhana dengan korelasi yang menurun pada skor tinggi, koreksi ini menaksir rendah korelasi dan $SD$ populasi (Tabel~\ref{tab:slodr_case3}). Taksiran pada $b = 0$ memulihkan nilai sebenarnya dengan baik. Estimasi $\bar{R} \approx 0,425$ dengan demikian lebih tepat dibaca sebagai batas bawah bila SLODR berlaku.

Implikasi kedua menyangkut translasi persentil pada Tabel~\ref{tab:ranking_komposit}, yang memakai satu korelasi untuk seluruh rentang skor. Bila korelasi menurun pada skor tinggi, simulasi menunjukkan bahwa ekor atas skor komposit menipis dan ekor bawah menebal (Tabel~\ref{tab:slodr_ekor}), sehingga kelangkaan skor ekstrem dapat keliru ditaksir pada kedua ujung. Arah efek ini berlawanan dengan efek marginal miring ke kanan pada Subbagian Implikasi terhadap Estimasi Persentil Nasional. Pada $b = 0,04$, skor komposit pada persentil normal ke-99,9 lebih rendah daripada pada korelasi konstan dan mendekati skor pada normal, sehingga kedua penyimpangan dari normal sebagian saling meniadakan di ekor atas. Hal ini berarti Tabel~\ref{tab:ranking_komposit} dibaca sebagai kisaran skenario, dan pengujian korelasi yang bergantung pada tingkat skor dengan data empiris menjadi arah pengembangan lanjutan (lihat Rekomendasi untuk Penelitian Selanjutnya).

\subsection{Struktur Faktor dan Keterbatasan Homogenitas Korelasi}
Hasil CFA (Tabel~\ref{tab:cfa_fit}--\ref{tab:factor_loadings}) memberi konfirmasi struktural atas gagasan \textit{g-factor} yang melatari ekspektasi korelasi positif antarsubtes pada Subbagian Interpretasi Nilai $r \approx 0,300$. Keunggulan Model 2 terhadap Model 1 menunjukkan bahwa kedelapan subtes UTBK-SNBT terbagi ke dalam dua domain yang secara teoretis dapat dibedakan, yaitu Literasi dan Numerasi. Meskipun demikian, korelasi antarfaktor yang tinggi pada CFA, dan yang lebih rendah pada EFA \textit{oblimin} yang mengizinkan \textit{cross-loading}, menunjukkan bahwa kedua domain tetap berbagi varians bersama yang besar. Pola itu konsisten dengan keberadaan \textit{g-factor} yang mendasari performa pada kedua domain sekaligus. Pola hierarkis ini selaras dengan \textit{Carroll's three-stratum theory} \autocite{carroll1993} (lihat Tinjauan Pustaka).

Temuan struktural ini turut menjelaskan pola deviasi pada diagnostik homogenitas korelasi (Tabel~\ref{tab:diagnostik}, Subbagian Evaluasi Homogenitas Struktur Korelasi Antarsubtes). Subtes dengan muatan faktor yang relatif lebih lemah pada CFA, seperti LBI Soshum, juga menunjukkan rasio $SD$ tersirat yang menyimpang cukup jauh dari $SD$ nasional. Hal ini mengindikasikan bahwa subtes itu berbagi lebih sedikit varians bersama dengan skor komposit dibandingkan subtes lain. Pola ini memperkuat argumen bahwa koreksi Case III, yang mengasumsikan korelasi seragam antarpasangan subtes, paling akurat diterapkan pada tingkat rerata populasi agregat (lihat Bagian Keterbatasan Penelitian).

\subsection{Implikasi terhadap Estimasi Persentil Nasional}
Atenuasi nilai $r \approx 0,300$ akibat \textit{restriction of range} membawa implikasi penting terhadap akurasi penerjemahan skor komposit menjadi peringkat persentil nasional (Tabel~\ref{tab:ranking_komposit}). Penggunaan korelasi \textit{raw sample} tanpa koreksi mempersempit estimasi $SD$ skor komposit populasi, sehingga secara sistematis menempatkan skor tinggi pada persentil yang lebih ekstrem (\textit{overoptimistic}) dibandingkan posisi riilnya. Pola ini paling menonjol pada titik \textit{benchmark} ekstrem seperti $M + 2\,SD$ hingga $M + 3\,SD$, tempat selisih skor komposit antarskenario korelasi mencapai puluhan poin (lihat Subbagian Matriks Translasi Persentil Standar).

Koreksi Case III mengoreksi distorsi itu pada tingkat agregat populasi. Namun, sebagaimana ditunjukkan pada Subbagian Struktur Faktor dan Keterbatasan Homogenitas Korelasi, koreksi ini tetap mengandung ketidakpastian residual akibat heterogenitas struktur korelasi antarsubtes yang riil. Akibatnya, translasi persentil pada Tabel~\ref{tab:ranking_komposit} paling tepat dibaca sebagai kisaran estimasi yang bergantung pada asumsi korelasi yang digunakan. Nuansa ini relevan bagi pembaca yang hendak memperkirakan posisi persentil nasionalnya sendiri.

Analisis ketidakpastian menambah satu pertimbangan praktis. Ketidakpastian pengambilan sampel pada $S_Z$ tergolong kecil, sebagaimana ditunjukkan oleh interval \textit{bootstrap} 95\% yang sempit pada $M + 2\,SD$ dan $M + 3\,SD$. Sebaliknya, ketergantungan pada $SD$ nasional tidak resmi jauh lebih besar. Apabila seluruh $SD$ nasional dikalikan $0,90$ hingga $1,10$, skor komposit pada $M + 2\,SD$ bergerak dari sekitar $608$ hingga $661$ (Subbagian Ketidakpastian Estimasi). Bagi pembaca yang ingin memperkirakan persentil nasionalnya dari skor komposit, rentang ini lebih informatif daripada interval \textit{bootstrap}, karena skor komposit yang bersesuaian dengan tiap persentil pada Tabel~\ref{tab:ranking_komposit} bergantung pada $S_Z$, yang pada gilirannya bergantung pada parameter nasional yang belum terverifikasi. Translasi skor mentah satu subtes ke persentil (Tabel~\ref{tab:ranking_subtes}) tidak bergantung pada korelasi antarsubtes, tetapi tetap bergantung pada $SD$ nasional subtes itu. Berdasarkan pertimbangan itu, hasil paper ini lebih tepat dipakai sebagai perkiraan kasar posisi relatif.

Asumsi distribusi normal menambah ketidakpastian yang tidak simetris. Tabel~\ref{tab:ranking_subtes} dan~\ref{tab:ranking_komposit} memakai persentil normal baku, padahal statistik nasional tidak resmi menunjukkan sebaran miring ke kanan dengan ekor kiri terpotong (Tabel~\ref{tab:nonnormal_diag}). Pada simulasi dengan marginal Johnson SB (Tabel~\ref{tab:nonnormal_sim}), hampiran normal relatif aman di sekitar $\pm 1\,SD$, tetapi tidak netral di ekor. Di ekor atas, kelompok peserta yang berskor $M + 2\,S_Z$ dan $M + 3\,S_Z$ sekitar $1,5$ dan $2,8$--$3$ kali lebih besar daripada yang diprediksi model normal, sehingga peringkat nasional yang diperkirakan dari persentil normal untuk skor tinggi cenderung terlalu optimistis. Arah ini sama dengan distorsi akibat penggunaan korelasi \textit{raw sample} yang dibahas di atas. Ukurannya, sekitar $13$--$15$ poin pada $M + 2\,S_Z$ dan $27$--$29$ poin pada $M + 3\,S_Z$, kira-kira setengah dari pergeseran skor akibat perubahan $SD$ nasional sebesar $\pm 10\%$ pada titik yang sama. Di ekor bawah, model normal menempatkan lebih banyak peserta daripada yang diindikasikan simulasi, tetapi ekor bawah ini kurang relevan secara praktis karena skor pada $M - 3\,SD$ sudah berada di bawah minimum nasional pada seluruh subtes (Subbagian Evaluasi Asumsi Distribusi). Karena bentuk distribusi skor komposit dihampiri dari tujuh statistik per subtes, yang dapat dinyatakan hanyalah bahwa asumsi normal tidak netral di ekor dan dampaknya tidak dapat diabaikan. Besaran koreksi yang tepat tidak dapat ditentukan tanpa data mikro nasional.